\section{Bibliography}

\subsection{Sources in Western Languages}

\bibentry{Anderson, Joseph L. and Donald Richie;\textit{The Japanese Film: Art and Industry};Princeton University Press;1982}
\bibentry{Betz, Betty;Be-bop, BG, and a Box of Bubble Gum;\textit{Down Beat;} 1949.8.26}
\bibentry{Boulton, David;\textit{Jazz in Britain};W.H. Allen, London;1958}
\bibentry{Brown, Sidney DeVere: original title: "Jazz in Japan," ;\textit{Japan: Ein Lesebuch}, tr. by Christian Schwandt, in Peter Portner, ed., Konkursbuch;University of Hamburg;1986:117-131頁}
\bibentry{Brown, Sidney DeVere;\textit{New Orleans in Tokyo: Nanri Fumio (1910-1975)-- Jazz Trumpet Virtuoso};Southwest Conference on Asian Studies, New Orleans;1989.11.3}
\bibentry{Brown, Sidney DeVere;\textit{Toshiko Akiyoshi: The Meeting of East and West in Jazz};Southwest Conference on Asian Studies, University of Texas, Austin;1990.10.13}
\bibentry{Collier, James Lincoln, \textit{The Making of Jazz: A Comprehensive History}, Houghton Mifflin, 1978}
\bibentry{Filipacchi, Daniel, Jazz Overseas: France, \textit{The Encyclopedia of Jazz}, ed. Leonard Feather, Horizon Press, 1960}
\bibentry{Franklin, A. David;A Preliminary Study of the Acceptance of Jazz by French Music Critics in the 1920s and Early 1930s, \textit{Annual Review of Jazz Studies}, 1985}
\bibentry{Hamm, Charles;\textit{Yesterdays: Popular Song in America};Norton \& Company, Inc.;1979}
\bibentry{Henthoff, Nat;Japan Getting Hip: Granz;\textit{Down Beat} 1954.2}
\bibentry{Hershey, Burnett;Jazz Latitude;\textit{New York Times}, 1922.6.25}
\bibentry{Kanahele, George S.;\textit{Hawaiian Music and Musicians, An Illustrated History};The University Press of Hawaii;1979}
\bibentry{Kater, Michael H.;Forbidden Fruit? Jazz in the Third Reich\textit{Journal of Modern History}, 1986}
\bibentry{Kennington;\textit{Internatinal Jazz Bibliography};1983}
\bibentry{Komiya Toyotaka, tr. by Edward G. Seidensticker and Donald Keene;\textit{Japanese Music and Drama in the Meiji Era};Obunsha;1956}
\bibentry{Kyogoku Takatoshi;Jazz with a classical tint rules Japan;\textit{Down Beat};1948.12.1}
\bibentry{Lapham, Claude;\textit{Looking at Japanese Jazz};Metronome, pp. 12,27;1936.6}
\bibentry{Lapham, Claude;Popular Japanese Music;\textit{Asia}, pg. 859;1937.12}
\bibentry{Lischeron, Beth;Jazz in Japan;\textit{Wingspan}「翼の王国」日空の機内誌 pg. 7;1989.10}
\bibentry{Malm, William P.; A Century of Proletarian Music In Japan;\textit{The Transactions of the Asiatic Society of Japan};1984.6}
\bibentry{Manuel, Peter;\textit{Popular Musics of the Non-Western World: An Introductory Survey};Oxford University Press;1988}
\bibentry{Martin, Harris I.;Popular Music and Social Change in Prewar Japan;\textit{The Japan Interpreter}, Vol. VII, Nos. 3-4,pg. 333;1972}
\bibentry{Muraoka Tay;Jazz Back in Japan;\textit{Metronome}, pg. 17;1947.8}
\bibentry{Nakazawa Mayumi;Japan Jazzline 1910-1970;\textit{Wingspan}「翼の王国」日空の機内誌 1989.10:7頁}
\bibentry{Ogren, Kathy J., \textit{The Jazz Revolution: Twenties America and the Meaning of Jazz},Oxford University Press, 1989}
\bibentry{Lyons, Len and Don Perlo, \textit{Jazz Portraits- The Lives and Music of the Jazz Masters}, Quill,1989}
\bibentry{Pyle, Kenneth B.;\textit{The New Generation in Meiji Japan};Problems of Cultural Identity, 1885-1895;Stanford University Press;1969}
\bibentry{Randal, Don;\textit{The New Harvard Dictionary of Music};Belknap Press;1986}
\bibentry{Richie, Donald;\textit{Japanese Cinema—An Introduction};Oxford University Press;1990}
\bibentry{Schuller, Gunther;\textit{Early Jazz: Its Roots and Musical Development};Oxford University Press;1968}
\bibentry{Schuller, Gunther;\textit{The Swing Era: The Development of Jazz 1930-1945};Oxford University Press;1989}
\bibentry{Seidensticker, Edward;\textit{Low City, High City};Charles E. Tuttle Co.:1983}
\bibentry{Seidensticker, Edward;\textit{Tokyo Rising};Alfred A. Knopf:1990}
\bibentry{Uenoda Setsuo;\textit{Japan and Jazz}教文館1930.11.28}
\bibentry{Waterhouse, David;Hogaku Preserved (I);\textit{Recorded Sound}, No. 33, pg.383-402;1965}
\bibentry{Waterhouse, David;Hogaku Preserved (II);\textit{Recorded Sound}, No. 42, pg.408-426;1974}
\bibentry{Waterhouse, David;Music: Western, \textit{Kodansha Encyclopedia of Japan}, Ed. Reischauer, Edwin O. and Tsuru Shigeto et al; Kodansha; 1983.11.25}
\bibentry{Jazz in China;\textit{Musical Courier};1924.6.26}
\bibentry{Jazz Grows Popular In China, Dance Record Sales Reveal;\textit{New York Times};1929.11.16}
\bibentry{Jazz in the Orient;\textit{New York Times};1929.11.19}
\bibentry{Music: Japanese Jazz Writers Get an American Radio Hearing;\textit{Newsweek};1937.6.26}
\bibentry{Jazzu;\textit{Time};1949.8.8}
\bibentry{Western Music in Japan;\textit{Music Educators Journal;} 1950.1}
\subsection{Sources in Japanese}

\bibentry{青木正、唐端勝、野川香文『軽音楽とそのレコード』三省堂 1938.10.5}
\bibentry{秋山邦晴『日本の映画音楽史1』田畑書店 1974.11.25}
\bibentry{浅香淳など編『新音楽辞典、楽語』音楽の友社1991.6.20}
\bibentry{池田恂一『昭和流行歌の軌跡』白馬出版 1985}
\bibentry{石川義弘『娯楽の戦前史』東書選書1981.12.14}
\bibentry{伊庭孝「民謡への公開状と抗議、軟弱・悪趣味の現代民謡」『読売新聞』1929.8.4}
\bibentry{岩村英武 american south. syncopated orchestra『詩と音楽』1923年6月号}
\bibentry{岩浪洋三『日本のジャズメン』大陸書房 1982.8.13}
\bibentry{内川芳美『現代史資料・マスメディア統制』みすず書房 1975.10.20}
\bibentry{内田晃一『日本のジャズ史』Swing Journal社1976.7.15}
\bibentry{海野弘「ジャズとモダン都市」『国文学』1990.2}
\bibentry{大森盛太郎『日本の洋楽』第一巻、新門出版社 1986.12.20}
\bibentry{永井良和「ダンスホールの百年」『流行の風俗学』、多田道太郎編 1986}
\bibentry{永井良和　連載「檻の中の舞踏会」『ダンスファン』1988.4から1989.6}
\bibentry{永井良和『社交ダンスと日本人』晶文社 1991.8.30}
\bibentry{長尾直『流行歌のイデオロギー』世界思想社1974.4.1}
\bibentry{加太こうじ『歌の昭和史』時事通信社 1975.9.20}
\bibentry{加太こうじ『流行歌論』東書選書1981.12.14}
\bibentry{兼常清佐「時代と流行歌曲の問題、日本の流行小唄はかなりに道徳的」『読売新聞』1929.8.16}
\bibentry{川路柳虹「流行歌曲の廃頽性(上)」『読売新聞』1929.8.10}
\bibentry{川路柳虹「時代と流行歌曲の問題、流行歌曲の廃頽性(下)」『読売新聞』1929.8.13}
\bibentry{北原鉄雄編「ジャズ音楽」『アルス音楽大講座第9巻』1936.4.20}
\bibentry{倉田喜弘『日本のレコード文化史』東京書籍 1979.3.28}
\bibentry{倉田喜弘『日本近代思想大系―芸能』岩波書店 1989.10.1}
\bibentry{小泉文夫『歌謡曲の構造』冬樹社 1984.11.20}
\bibentry{小泉文夫『民族音楽の世界』日本放送出版協会 1985.7.20}
\bibentry{小宮豊隆『明治文化史』洋々社 1954.6.1}
\bibentry{西條八十「民謡への公開状と抗議、伊庭孝氏に與ふ―東京行進曲と僕―」『読売新聞』1929.8.4}
\bibentry{斉藤龍『横浜音楽活動史の研究ー大正期の大衆音楽ー』鶴見大学紀要第24号第3部・保育・保健歯科編、別刷、1987.3}
\bibentry{斉藤憐『昭和のバンスキングたち』ミュージック・マガジンの本 1983.3.20}
\bibentry{坂本圭太郎『物語・軍歌史』創思社 1984.4.25}
\bibentry{塩入亀輔『ジャズ音楽』敬文館 1929.9.5}
\bibentry{塩入亀輔『ジャズ音楽ー普及版』敬文館 1934.1.10}
\bibentry{塩入亀輔『ジャズ音楽入門』音楽世界社 1936.8.10}
\bibentry{塩澤実信『昭和のすたるじい流行歌』第三文明社 1991.8.15}
\bibentry{篠崎正『ジャズ鑑賞の仕方』金園社 1955.7.20}
\bibentry{神保環一郎『レコード音楽第一課』創元社 1935.12.23}
\bibentry{瀬川昌久「米国におけるスイングへのノスタルジア」『Jazz』1974.11}
\bibentry{瀬川昌久「音楽への愛着」(連載)『Jazz』75年3月-76年12月}
\bibentry{瀬川昌久『ジャズで踊って』サイマル出版会 1983.8}
\bibentry{園部三郎『日本民衆歌謡史考』朝日新聞社 1980.8.20}
\bibentry{園部三郎『日本の流行歌、その魅力と流行のしくみ』大月書店 1980.6.23}
\bibentry{園部三郎『音楽五十年史』時事通信社 1956.3.10}
\bibentry{田中純一郎『日本映画発達史』中央公論社 1968.7.25}
\bibentry{千葉亀雄「新感覚論」『文芸日本』1925.4}
\bibentry{辻厚生、西山卓『昨日・今日・明日に歌う』飛鳥書房、奈良市 1985.7.7}
\bibentry{中澤まゆみ「日本人がジャズを口ずさんだ日」『潮』1981.12}
\bibentry{中野重治「時代と流行歌曲の問題、も一つの東京行進曲へ」『読売新聞』1929.8.20}
\bibentry{中山卯郎『中山晋平作曲目録譜』豆々樹社 1980.2.29}
\bibentry{夏川太郎訳『ジャヅ・ホワイトマン著』田舎社;1929.9.1 (原著\textit{Jazz}, by Paul Whiteman and M. McBridge, 1926)}
\bibentry{成田龍一「都市の生活革命」大正演劇研究会編『大正の演劇と都市』武蔵野書房19914.6}
\bibentry{西沢爽『雑学歌謡昭和史』毎日新聞社 1980.10.30}
\bibentry{西沢爽『日本近代歌謡史』(上下、付録、別冊付録) 桜楓社 1990.11.20}
\bibentry{野川香文『日本のジャズ演奏家』音楽芸術、1951.11}
\bibentry{野口久光訳『ジャズへの道』Ulanoff, Barry著、新興楽譜出版社;1955,10,1}
\bibentry{服部龍太郎『ジャズ楽器の奏法』春陽堂 1931.5.15}
\bibentry{服部龍太郎『ライター・ミュジック』音楽の友社 1957.6.30}
\bibentry{服部良一『ぼくの音楽人生』中央交芸社 1982.10.1}
\bibentry{早津敏彦『アロハ！メレ・ハワイ、日本ハワイ音楽・舞踊史』サンクリエイト 1986}
\bibentry{平岡正明『大歌謡論』筑摩書房1989.8.30}
\bibentry{福田俊二編『日本流行歌年表』彩工社 1968.6.15}
\bibentry{二村定一訳『ジャズ音楽』大衆公論社;1930.11.10(原著 \textit{Jazz}, by Paul Whiteman and M. McBridge, 1926)}
\bibentry{茂田信男、矢沢保、島田芳文、横沢千秋『日本流行歌史』社会思想社 1981.1.30}
\bibentry{細川周平　西洋音楽の日本化・大衆化『Music Magazine』1989.4連載開始}
\bibentry{堀内敬三「時代と流行歌曲の問題、流行歌は進歩する(上)」『読売新聞』1929.8.14}
\bibentry{堀内敬三「時代と流行歌曲の問題、流行歌は進歩する(下)」『読売新聞』1929.8.15}
\bibentry{堀内敬三編『世界音楽全集』「第十九巻、流行歌曲全集」春秋社 1931.4.15}
\bibentry{堀内敬三『音楽五十年史』鱒書房1942.12.10}
\bibentry{堀内敬三「ジャズ鳥瞰」『トオキイ音楽』1934.9}
\bibentry{堀内敬三『データ・音楽・日本』民主音楽協会 1955.2.20}
\bibentry{室伏高信『街頭の社会学』田舎社 1929.3.5}
\bibentry{清水正一『ジャズの歴史』創元新社 1957.6.10}
\bibentry{清水正一『ディキシーランドジャズ入門』荒地出版社 1967.10.15}
\subsection{Special Issues and Features}

\bibentry{「日本にジャズが入ってきた頃 （創始期から1960年まで）」『季刊ジャズ批評』 1972.6.20}
\bibentry{「ジャズ批評 創刊号55年版」『季刊ジャズ批評』1980.4.20}
\bibentry{「日本のジャズ」（別冊1億人の昭和史）毎日新聞社 1982}
\bibentry{「ジャズ批評 創刊号61年版」『季刊ジャズ批評』1986.5.20}
\bibentry{「日本音楽著作権協会50年史」日本音楽著作権協会 1989.11.18}
